\documentclass[10pt,a4paper]{article}
\usepackage[margin=20mm,headheight=15pt,headsep=7mm,footskip=10mm]{geometry}
\usepackage{fontspec,amsmath,amssymb,mathtools,microtype}
\setmainfont{lmroman10-regular.otf}[BoldFont=lmroman10-bold.otf,ItalicFont=lmroman10-italic.otf,BoldItalicFont=lmroman10-bolditalic.otf]
\setsansfont{lmsans10-regular.otf}[BoldFont=lmsans10-bold.otf,ItalicFont=lmsans10-oblique.otf,BoldItalicFont=lmsans10-boldoblique.otf]
\setmonofont[Scale=.86]{lmmono10-regular.otf}[ItalicFont=lmmono10-italic.otf]
\usepackage{ragged2e}\usepackage{xcolor,booktabs,longtable,array,graphicx,enumitem,fancyhdr,needspace,titlesec,hyperref,xurl}
\definecolor{AtlasNavy}{HTML}{18334C}\definecolor{AtlasTeal}{HTML}{116D73}\definecolor{AtlasGray}{HTML}{52616D}
\hypersetup{colorlinks=true,linkcolor=AtlasTeal,urlcolor=AtlasTeal,pdfcreator={Complexity Atlas},unicode=true}
\urlstyle{same}\setlength{\parindent}{0pt}\setlength{\parskip}{4.8pt}\setlength{\emergencystretch}{3em}
\renewcommand{\baselinestretch}{1.08}\setcounter{secnumdepth}{0}\setlist{nosep,leftmargin=1.5em}
\titleformat{\section}{\Large\bfseries\sffamily\color{AtlasNavy}\raggedright}{}{0em}{}
\titleformat{\subsection}{\large\bfseries\sffamily\color{AtlasTeal}\raggedright}{}{0em}{}
\titleformat{\subsubsection}{\normalsize\bfseries\color{AtlasNavy}\raggedright}{}{0em}{}
\titlespacing*{\section}{0pt}{14pt}{6pt}\titlespacing*{\subsection}{0pt}{10pt}{5pt}
\pagestyle{fancy}\fancyhf{}\fancyhead[L]{\small\sffamily\color{AtlasGray}COMPLEXITY ATLAS}
\fancyhead[R]{\small\sffamily\color{AtlasGray}ENGLISH EDITION}\fancyfoot[C]{\small\thepage}\renewcommand{\headrulewidth}{.3pt}
\newcommand{\AtlasMath}[1]{\par\begingroup\centering\setbox0=\hbox{$\displaystyle #1$}\ifdim\wd0>1.4\linewidth\typeout{ATLAS WIDE FORMULA: \the\wd0}\fi\ifdim\wd0>\linewidth\resizebox{\linewidth}{!}{\box0}\else\box0\fi\par\endgroup}
\newcommand{\AtlasInlineMath}[1]{\begingroup\setbox0=\hbox{$#1$}\ifdim\wd0>.4\linewidth\endgroup\AtlasMath{#1}\else\endgroup$#1$\fi}
\pdfstringdefDisableCommands{\def\AtlasInlineMath#1{#1}}
\DeclareMathOperator{\polylog}{polylog}\DeclareMathOperator{\dist}{dist}\DeclareMathOperator{\diam}{diam}
\DeclareMathOperator{\rank}{rank}\DeclareMathOperator{\tr}{tr}\DeclareMathOperator{\diag}{diag}
\DeclareMathOperator{\prox}{prox}\DeclareMathOperator{\KL}{KL}\DeclareMathOperator{\Var}{Var}
\providecommand{\eps}{\varepsilon}\providecommand{\R}{\mathbb R}\providecommand{\E}{\mathbb E}
\providecommand{\mathbbm}[1]{\mathbf{#1}}\providecommand{\mathds}[1]{\mathbf{#1}}
\providecommand{\llbracket}{[\![}\providecommand{\rrbracket}{]\!]}
\begin{document}
\raggedright

{\LARGE\bfseries\sffamily\color{AtlasNavy} Strongly monotone variational inequalities: fresh operator samples and mean-square distance\par}\vspace{5pt}
\textbf{Setting ID.} \texttt{CV\_\allowbreak{}S\_\allowbreak{}STRONG\_\allowbreak{}VI\_\allowbreak{}FRESH}

Variational inequalities

Solve \AtlasInlineMath{0\in F(z)+N_K(z)}.

\textbf{Status.} Bounds missing

\textbf{TASK tags.} solution concept: Solution distance; task family: Variational inequality

\textbf{MODEL tags.} initialization: Initial distance; initialization: Initialization; monotonicity: Strong monotonicity; monotonicity: Strongly monotone

\textbf{ORACLE tags.} feedback quality: Exact; feedback quality: Stochastic unbiased bounded variance; interface family: Fresh operator; operation: Projection; operation: Stochastic operator value; public information: Hidden objects explicitly declared; public information: Public inputs explicitly declared

\textbf{ALGORITHMS tags.} adaptivity: Adaptive; randomization: Private randomness allowed; restriction: Explicitly no additional algorithm restrictions; restriction: No additional span restriction explicitly imposed; restriction: No additional zero-respecting restriction explicitly imposed

\textbf{OUTPUT tags.} certificate: No additional output certificate required; object: Point

\textbf{METRIC tags.} criterion: Distance

\textbf{GUARANTEE tags.} aggregation: Per instance; aggregation: Second moment of the complete error; mode: Expectation; quantifier: Algorithm before instance; randomness: Algorithm private seed; randomness: Oracle samples

\textbf{RESOURCE tags.} accounting: Joint coordinates retained; budget: Fixed worst-case budget; measure: N operator sample; measure: N projection; setup: Initialization included

\needspace{7\baselineskip}\section{1. Setting specification}
\needspace{4\baselineskip}\subsection{Task}
Solve \AtlasInlineMath{0\in F(z)+N_K(z)}.

\needspace{4\baselineskip}\subsection{Admissible model}
The set \AtlasInlineMath{K\subset\mathbb{R}^d} is known, nonempty, closed, and convex. The operator \AtlasInlineMath{F} is \AtlasInlineMath{\beta}-Lipschitz and \AtlasInlineMath{\mu}-strongly monotone: \AtlasInlineMath{\langle F(u)-F(v),u-v\rangle\ge\mu\|u-v\|_2^2} and \AtlasInlineMath{\beta\ge\mu>0}.

A unique solution \AtlasInlineMath{z^\star} exists, and the public initial point \AtlasInlineMath{z_0\in K} satisfies \AtlasInlineMath{\|z_0-z^\star\|_2\le R_0}.

\needspace{4\baselineskip}\subsection{Oracle interface}
Public inputs include dimension \AtlasInlineMath{d} and \AtlasInlineMath{R_0}, \AtlasInlineMath{\beta}, \AtlasInlineMath{\varepsilon}, \AtlasInlineMath{\mu}, and \AtlasInlineMath{\sigma}, together with the stated initial points and feasible-set information. Analytic representations of hidden functions, operators, and matrices are not supplied.


\AtlasMath{z\in K\longmapsto G_F(z;\xi)}


Let \AtlasInlineMath{\mathcal H_t} denote the full history before query t. Conditional on \AtlasInlineMath{\mathcal H_t,z_t}, the response \AtlasInlineMath{G_t\in\mathbb R^d} satisfies \AtlasInlineMath{\mathbb E[G_t\mid\mathcal H_t,z_t]=F(z_t)} and \AtlasInlineMath{\mathbb E[\|G_t-F(z_t)\|_2^2\mid\mathcal H_t,z_t]\le\sigma^2}. Each call uses an independent fresh sample.

No reusable sample identifier or cross-point sample reuse is available. No mean-square smoothness of sample gradients or operators is additionally assumed. Blocks within a full vector may be correlated.

Exact Euclidean projection onto the known set is allowed. Its calls are recorded separately and excluded from the primary query count.

\needspace{4\baselineskip}\subsection{Output}
Return \AtlasInlineMath{\widehat{z}\in K}.

The output may be any measurable postprocessing of the history, including an unqueried point, average, or convex combination. It must satisfy the stated domain and feasibility requirements; no additional function value or computable certificate is required.

\needspace{4\baselineskip}\subsection{Loss or accuracy criterion}

\AtlasMath{\|\widehat{z}-z^\star\|_2}


\needspace{4\baselineskip}\subsection{Resource accounting}
One operator-value sample costs one \AtlasInlineMath{N_{\mathrm{sample}}} call. Projection count is stored separately.

Use a fixed worst-case budget: calls never exceed the prespecified \AtlasInlineMath{N} on any instance or random execution path. Earlier stopping within the budget is permitted.

Count all information calls for one instance and one requested accuracy, including initialization. Exact-real arithmetic, storage, and solving known models from public inputs and returned data are uncharged. This query-complexity quantity does not certify arithmetic or bit complexity.

\needspace{4\baselineskip}\subsection{Algorithm views and success guarantees}
Arbitrary measurable sequential adaptive randomized algorithms are allowed. The private seed is independent of the hidden instance. Queries, stopping, and output may use the full history, without additional span, zero-respecting, communication, or memory restrictions.

Choose one algorithm, then require the guarantee for every admissible instance and stochastic-feedback mechanism.


\AtlasMath{\mathbb{E}\!\left[\left(\|\widehat{z}-z^\star\|_2\right)^2\right]\le\varepsilon^2.}


Expectation is over private algorithmic randomness and the random samples actually used.

\textbf{Budget convention.} fixed\_worst\_case\_cap

\needspace{7\baselineskip}\section{2. Notation and parameters}
\AtlasInlineMath{R_0}: Euclidean distance bound from the public initial point to the same solution promised by the model; not the entire domain diameter \textbf{Domain:} A nonnegative real.

\AtlasInlineMath{\|z_0-z^\star\|_2\le R_0}

\AtlasInlineMath{\beta}: Euclidean Lipschitz upper bound for the population operator on the declared query domain \textbf{Domain:} Positive real number.

\AtlasInlineMath{\|F(z)-F(w)\|_2\le\beta\|z-w\|_2}

\AtlasInlineMath{d}: Finite Euclidean dimension of the operator domain; regime quantifiers specify whether complexity is uniform in dimension \textbf{Domain:} Positive integer.

\AtlasInlineMath{\varepsilon}: Euclidean distance tolerance to the unique solution; the mean-square contract targets squared distance \AtlasInlineMath{\varepsilon^2}. \textbf{Domain:} Positive real number.

\AtlasInlineMath{\mu}: A strong-monotonicity lower bound for the full operator. \textbf{Domain:} Positive real number.

\AtlasInlineMath{\langle F(u)-F(v),u-v\rangle\ge\mu\|u-v\|_2^2}

\AtlasInlineMath{\sigma}: Standard-deviation upper bound for noise in the full stochastic response, controlling the joint squared norm rather than summing independent block budgets. \textbf{Domain:} A nonnegative real.

\AtlasInlineMath{\mathbb E[\|G_t-\mathbb E[G_t\mid\mathcal H_t,z_t]\|_2^2\mid\mathcal H_t,z_t]\le\sigma^2}

\AtlasInlineMath{\|\cdot\|_2,\langle\cdot,\cdot\rangle}: Euclidean norm and inner product; the product-space norm is the square root of the sum of the squared block norms. \textbf{Domain:} The specified finite-dimensional real vector space.

\AtlasInlineMath{\|v\|_2=\sqrt{\langle v,v\rangle}}

\AtlasInlineMath{F}: The unknown operator is single-valued. Monotonicity, strong monotonicity, or weak Minty properties use the specification’s explicit inequalities. \textbf{Domain:} A vector map from the declared query domain into \AtlasInlineMath{\mathbb R^d}.

\AtlasInlineMath{z,u,v}: The operator argument and comparison point in the inequality. \textbf{Domain:} The ambient space is \AtlasInlineMath{\mathbb R^d}. Formulas constrained by \AtlasInlineMath{K} require the corresponding point to lie in \AtlasInlineMath{K}.

\AtlasInlineMath{K}: A known nonempty closed convex feasible set, possibly unbounded. \textbf{Domain:} A subset of Euclidean space.

\AtlasInlineMath{K\subset\mathbb R^d}

\AtlasInlineMath{z_0,z^\star}: Public initial point and the same solution appearing in the model assumption \textbf{Domain:} Feasible domain of this profile; the distance bound refers to the same solution.

\AtlasInlineMath{\widehat z}: The point or variable blocks ultimately returned by the algorithm, corresponding to the hatted variables in the error formula. \textbf{Domain:} \AtlasInlineMath{K\subset\mathbb R^d}.

\AtlasInlineMath{G_F(z;\xi),G_t}: Random vector sample of population operator \AtlasInlineMath{F} and its response at time \AtlasInlineMath{t}. \textbf{Domain:} \AtlasInlineMath{\mathbb R^d}.

\AtlasInlineMath{\xi,\mathcal H_t,z_t}: Current fresh-sample seed, complete pre-query history, and query point. \textbf{Domain:} The seed space, history-generated \AtlasInlineMath{\sigma}-algebra, and declared query domain.

\AtlasInlineMath{\mathbb E,\Pr}: Expectation and probability over the explicitly allowed algorithm and sample randomness, with the hidden instance fixed. \textbf{Domain:} Operations on the corresponding probability space.

\AtlasInlineMath{N_{\rm sample}}: Number of fresh-sample complete-operator-vector calls \textbf{Domain:} Nonnegative integer count; expected work is its expectation.

\AtlasInlineMath{N_{\rm projection}}: Projection count, uncharged in the primary query-cost projection \textbf{Domain:} Nonnegative integer count; expected work is its expectation.

\AtlasInlineMath{N}: A uniform cap on principal information calls for every admissible instance and execution path. \textbf{Domain:} Nonnegative integer.

\AtlasInlineMath{N^\star}: The infimum principal query budget among algorithms satisfying every stated instancewise output guarantee. \textbf{Domain:} Nonnegative extended real value; by convention, positive infinity if no successful algorithm or no uniform finite budget exists.

\AtlasInlineMath{\mathcal A,I}: An algorithm selected from public inputs, a hidden admissible object fixed before interaction, and an oracle law satisfying the declared conditional means, variances, and independence protocol. \textbf{Domain:} The specified algorithm and model classes.

\AtlasInlineMath{\omega}: A complete realization of algorithmic and sample randomness. \textbf{Domain:} Corresponding sample space.

\AtlasInlineMath{O,\Omega,\Theta,\widetilde O}: Asymptotic upper and lower bounds, constant-factor matching, and upper bounds hiding the stated logarithmic factors, within the corresponding regions. \textbf{Domain:} Nonnegative complexity function; subscript p permits additional dependence only on fixed p.

\AtlasInlineMath{\log,\ln,\log_+}: Natural logarithm and its nonnegative truncation \textbf{Domain:} Arguments of log and ln are positive; the argument of \AtlasInlineMath{\log_+} may be zero.

\AtlasInlineMath{\log_+(t)=\log(\max\{1,t\})}

\AtlasInlineMath{\mathfrak A_\varepsilon}: Algorithms using the allowed interfaces, algorithm restrictions, and output domain, and satisfying the displayed complete error guarantee for every fixed admissible instance \textbf{Domain:} The algorithm class is defined at fixed public parameters and target tolerance \AtlasInlineMath{\varepsilon}.

\needspace{12\baselineskip}\section{3. Complexity bracket}
\needspace{4\baselineskip}\subsection{Compared quantities}
\needspace{4\baselineskip}\subsubsection{Q1 - target accuracy cost}
\AtlasMath{N^\star(\varepsilon)=\inf_{\mathcal A\in\mathfrak A_\varepsilon}\sup_I \sup_\omega N_{\rm sample}(\mathcal A,I,\omega)}
\textbf{Scope.} The instancewise guarantee uses the same contract; dimension obeys the specification and theorem region.

\textbf{Optimization direction.} minimize

One call is charged for the full response specified by the interface. The algorithm class obeys the declared randomization, adaptivity, and output rules and must satisfy the complete error guarantee on every fixed admissible instance. The call cap is uniform over instances and execution paths.

\needspace{4\baselineskip}\subsection{R1 - The nontrivial accuracy region.}
\textbf{Region.} \AtlasInlineMath{0<\varepsilon<R_0}.

\textbf{Quantity.} Q1

\textbf{Resource.} The principal resource specified in the resource section. Primitive operations and unit costs are measured along this axis.

\textbf{Status.} Lower bound missing

\textbf{Lower bound.} No admitted bound for this comparison

\textbf{Upper bound.} B1

\AtlasMath{O\!\left(\frac{\beta}{\mu}\left[1+\log\frac{R_0}{\varepsilon}\right]+\frac{\sigma^2}{\mu^2\varepsilon^2}\right)}
\textbf{Comparison argument.} A matching lower bound for arbitrary randomized or deterministic algorithms is not recorded on this full parameter range. Noise-dominated subregions are listed separately.

\needspace{4\baselineskip}\subsection{R2 - Zero-query regime}
\textbf{Region.} \AtlasInlineMath{\varepsilon\ge R_0}.

\textbf{Quantity.} Q1

\textbf{Resource.} The principal resource specified in the resource section. Primitive operations and unit costs are measured along this axis.

\textbf{Status.} Optimal

\textbf{Lower bound.} Nonnegativity of the compared quantity

\AtlasMath{0}
\textbf{Upper bound.} B2

\AtlasMath{0}
\textbf{Comparison argument.} Cost is nonnegative.

\needspace{4\baselineskip}\subsection{R3 - Noise-dominated regime}
\textbf{Region.} \AtlasInlineMath{0<\varepsilon\le R_0/4}, \AtlasInlineMath{\sigma^2/(\mu^2\varepsilon^2)\ge1+(\beta/\mu)[1+\log(R_0/\varepsilon)]}.

\textbf{Quantity.} Q1

\textbf{Resource.} The principal resource specified in the resource section. Primitive operations and unit costs are measured along this axis.

\textbf{Status.} Optimal

\textbf{Lower bound.} B3

\AtlasMath{\Omega(\sigma^2/(\mu^2\varepsilon^2))}
\textbf{Upper bound.} B1

\AtlasMath{O(\sigma^2/(\mu^2\varepsilon^2))}
\textbf{Gap or comparison.} \AtlasInlineMath{O(1)}

\textbf{Comparison argument.} Matching is certified only in this noise-dominated subregion. The statistical lower bound is not a deterministic condition-number lower bound.

\needspace{7\baselineskip}\section{4. Bounds and supporting evidence}
The following conditions apply to every bound in this profile, in addition to each bound's stated region and algorithm scope.

\textbf{Dimension conditions.} Subject to the stated regime predicate; optimality in fixed small dimension is not asserted.

\textbf{Oracle mapping.} All primitive information operations are charged; see the specification and method description.

\textbf{Guarantee mapping.} The guarantee matches the specification. Any reduction is recorded in conditions.

\textbf{Resource mapping.} Complete query budget of a single algorithm

\needspace{4\baselineskip}\subsection{B1 - Geometrically restarted extragradient method (Corollary 1)}
\textbf{Admission.} Included in the bracket

\textbf{Direction.} upper

\textbf{Witness or lower-bound scope.} Randomized witness

\textbf{Algorithm-scope qualification.} The restarted method for strongly monotone stochastic operators uses fresh samples; randomness comes from the oracle.

\AtlasMath{O\!\left(\frac{\beta}{\mu}\left[1+\log\frac{R_0}{\varepsilon}\right]+\frac{\sigma^2}{\mu^2\varepsilon^2}\right)}
\textbf{Sources.} \href{https://arxiv.org/pdf/2609.04188v2}{[1]}

\textbf{Conditions.} Set squared-distance tolerance in Corollary 1 to \AtlasInlineMath{\delta=\varepsilon^2}, with \AtlasInlineMath{D_0=R_0} and \AtlasInlineMath{L=\beta}. Equations (2.43)-(2.46) give the mean-square distance bound with fixed stage count and batch sizes.

\textbf{Constants.} Constants are universal except for the stated fixed structural parameters; unknown thresholds retain symbolic values.

\textbf{Source location.} \href{https://arxiv.org/pdf/2609.04188v2}{[1]} theorem: Corollary 1; section: Section 2.3; pdf pages: 13; 14; equations: (2.43); (2.44); (2.45); (2.46)

\AtlasMath{O\!\left(\frac{\beta}{\mu}\left[1+\log\frac{R_0}{\varepsilon}\right]+\frac{\sigma^2}{\mu^2\varepsilon^2}\right)}
Set squared-distance tolerance in Corollary 1 to \AtlasInlineMath{\delta=\varepsilon^2}, with \AtlasInlineMath{D_0=R_0} and \AtlasInlineMath{L=\beta}. Equations (2.43)-(2.46) give the mean-square distance bound with fixed stage count and batch sizes.

\needspace{4\baselineskip}\subsection{B2 - Return the public initial point.}
\textbf{Admission.} Included in the bracket

\textbf{Direction.} upper

\textbf{Witness or lower-bound scope.} Deterministic witness

\textbf{Algorithm-scope qualification.} In this region, the witness returns the public initial point or known root, deterministically and without querying.

\AtlasMath{0}
\textbf{Conditions.} This follows directly from the initial-distance assumption.

\textbf{Constants.} Constants are universal except for the stated fixed structural parameters; unknown thresholds retain symbolic values.

\textbf{Local argument.} A1

\needspace{4\baselineskip}\subsection{B3 - Two-point Gaussian test}
\textbf{Admission.} Included in the bracket

\textbf{Direction.} lower

\textbf{Witness or lower-bound scope.} All randomized adaptive algorithms

\textbf{Algorithm-scope qualification.} Two-point testing, zero-query indistinguishability, and nonexistence counterexamples cover arbitrary privately randomized policies. These are not linear-span lower bounds.

\AtlasMath{\Omega(\sigma^2/(\mu^2\varepsilon^2))}
\textbf{Conditions.} In one dimension use \AtlasInlineMath{F_a(z)=\mu(z-a)}, \AtlasInlineMath{a\in\{-4\varepsilon,4\varepsilon\}}, and additive \AtlasInlineMath{N(0,\sigma^2)} noise. Subtract the known \AtlasInlineMath{\mu z} to obtain observations with mean \AtlasInlineMath{-\mu a}. Mean-square distance \AtlasInlineMath{\le\varepsilon^2} forces sign-testing error \AtlasInlineMath{\le1/16}. The per-sample KL divergence \AtlasInlineMath{32\mu^2\varepsilon^2/\sigma^2} then gives the stated two-point lower bound.

\textbf{Constants.} Constants are universal except for the stated fixed structural parameters; unknown thresholds retain symbolic values.

\textbf{Local argument.} A2

\needspace{7\baselineskip}\section{5. Local arguments}
\needspace{4\baselineskip}\subsection{A1 - Return the public initial point.}
\textbf{Scope.} All model, oracle, algorithm, guarantee, quantity, resource-projection, and region conditions attached to the mapped bounds remain in force.

The public initialization promise gives \AtlasInlineMath{\|z_0-z^\star\|\le R_0}. When \AtlasInlineMath{\varepsilon\ge R_0}, returning \AtlasInlineMath{z_0} without querying satisfies both the distance target and its mean-square version, deterministically. Query costs are nonnegative, so the minimum is exactly zero.

\needspace{4\baselineskip}\subsection{A2 - Gaussian two-point testing}
\textbf{Scope.} All model, oracle, algorithm, guarantee, quantity, resource-projection, and region conditions attached to the mapped bounds remain in force.

In one dimension use \AtlasInlineMath{F_a(z)=\mu(z-a)}, \AtlasInlineMath{a\in\{-4\varepsilon,4\varepsilon\}}, and additive \AtlasInlineMath{N(0,\sigma^2)} noise. Subtract the known \AtlasInlineMath{\mu z} to obtain observations with mean \AtlasInlineMath{-\mu a}. Mean-square distance \AtlasInlineMath{\le\varepsilon^2} forces sign-testing error \AtlasInlineMath{\le1/16}. The per-sample KL divergence \AtlasInlineMath{32\mu^2\varepsilon^2/\sigma^2} then gives the stated two-point lower bound.

\needspace{7\baselineskip}\section{References}
[1] Yao Ji; Guanghui Lan; Jason Zhu. Computation of Strong Solutions to Stochastic Variational Inequalities. 2026. arXiv:2609.04188v2, 2026-09-17.

\url{https://arxiv.org/pdf/2609.04188v2}

\end{document}
